\section{Related work}
\subsection{Aerial interaction and embodied evaluation}
AerialVLN, CityNav, and OpenFly emphasize language-conditioned aerial navigation, while AeroVerse and BEDI broaden evaluation toward perception, spatial reasoning, planning, motion, and perception--decision--action loops~\cite{liu2023aerialvln,lee2025citynav,gao2026openfly,yao2024aeroverse,guo2026bedi}. CityEQA evaluates question answering after active exploration in a realistic urban simulator~\cite{zhao2025cityeqa}. These systems support action-dependent observations and large-scale evaluation of aerial embodied behavior. Several also provide substantially richer 3D simulation, real-world data, or broader embodied tasks than the controlled orthographic setting used here.

\subsection{Disaster and remote-sensing agents}
xBD supplies paired disaster imagery and building-damage annotations~\cite{gupta2019xbd}. RescueADI, ThinkGeo, and DisasterM3 study tool use or multimodal disaster interpretation from remote-sensing records~\cite{liu2024rescueadi,thinkgeo2025,wang2025disasterm3}. These resources motivate disaster-specific questions. Evaluating an answer from a static image, however, does not measure the value of an agent's decision to observe again. DisasterClaw binds question scope, vehicle state, rendered observation, and terminal answer in one record. This binding also exposes a failure mode: a view that improves one building prediction can remove evidence required by a count or spatial question.

Table~\ref{tab:capability-comparison} summarizes the evaluation capabilities most relevant to this distinction. The comparison focuses on evaluation capabilities; it does not rank the overall scope of the benchmarks. Among the compared benchmarks, DisasterClaw explicitly standardizes a fixed annotation-derived geographic task and answer under view change, geographic evidence traces across observations, and a matched fixed-view answer call. These features define its contribution within a controlled observation setting.

\begin{table*}[!t]
\centering
\scriptsize
\caption{Capability-oriented comparison of benchmarks most relevant to task-dependent observation value.}
\label{tab:capability-comparison}
\resizebox{\textwidth}{!}{%
\begin{tabular}{@{}llccccc@{}}
\toprule
Benchmark & Primary task & \shortstack{Active\\view change} & \shortstack{Disaster-\\specific} & \shortstack{Fixed geographic\\task/answer\\under view change} & \shortstack{Answer after\\active sensing} & \shortstack{Explicit observation-\\value isolation}\\
\midrule
AerialVLN & Aerial vision--language navigation & $\checkmark$ & --- & --- & --- & ---\\
OpenFly & Large-scale aerial vision--language navigation & $\checkmark$ & --- & --- & --- & ---\\
BEDI & UAV embodied-agent evaluation & $\checkmark$ & $\triangle$ & --- & $\triangle$ & ---\\
CityEQA & Embodied question answering & $\checkmark$ & --- & $\triangle$ & $\checkmark$ & ---\\
RescueADI & Disaster remote-sensing agent & --- & $\checkmark$ & --- & $\triangle$ & ---\\
ThinkGeo & Tool-augmented remote-sensing agent & $\triangle$ & $\triangle$ & --- & $\triangle$ & ---\\
DisasterM3 & Disaster vision--language benchmark & --- & $\checkmark$ & --- & --- & ---\\
\textbf{DisasterClaw} & Disaster inspection and observation-value measurement & $\checkmark$ & $\checkmark$ & $\checkmark$ & $\checkmark$ & $\checkmark$\\
\bottomrule
\end{tabular}}
\vspace{2pt}
\parbox{\textwidth}{\footnotesize ``---'' means that the capability is not the primary evaluation contract reported by the work, rather than that it is technically impossible. $\triangle$ denotes a related but non-equivalent capability. DisasterClaw changes controlled orthographic image footprints rather than physical 3D viewpoints.}
\end{table*}


\subsection{Task value of an observation}
Active perception treats sensing as a task-directed action~\cite{bajcsy1988active}. Entropy and calibrated probabilities can guide such actions, but neither directly measures the correctness of a later task answer~\cite{guo2017calibration,ovadia2019can}. We therefore report localization, answer corrections and harms, executed rechecks, and terminal accuracy separately. The task-conditioned rule in this paper is a fixed heuristic for descending subject to predicted ROI coverage and motion cost. Complete-policy comparisons measure outcomes at the realized operating points. We additionally use fixed-view matched calls at executed recheck transitions to control for the extra answer opportunity and examine the contribution of changed observations and evidence within this diagnostic design.
